% =============================================================================
% 3. METHODOLOGY
% =============================================================================
\section{Methodology}
\label{sec:methodology}

\subsection{Research design}
\label{subsec:design}

The study adopts a longitudinal single-case study design
\citep{Patel2023,ORreilly2024} conducted over twelve calendar months at a
multinational aircraft-engine \MRO{} facility certified under \AS{} and
operating under regulatory oversight of an \ICAO{}-aligned civil aviation
authority. The case-study approach is appropriate because the research
questions are explanatory (\textit{how} and \textit{what changes}), the
unit of analysis (the training process) is contemporary, embedded in its
operational context, and not controllable by the researchers.

The intervention consists of the operationalisation of the Kirkpatrick
four-level model within the existing training process. The study compares
operational outcomes \textit{before} the intervention (baseline, derived from
the facility's existing quality records for the preceding twelve months) and
\textit{after} the intervention (intervention year), holding workforce size
and engine programmes constant to the extent feasible. The redesigned process
is presented in Section~\ref{subsec:results-process}.

\subsection{Data collection}
\label{subsec:data}

Data collection combined four sources to support triangulation:

\begin{enumerate}[label=(\alph*)]
  \item \textbf{Post-training questionnaires} (Level~1 and Level~2): a
        seven-item satisfaction questionnaire administered immediately after
        every training session (n~=~10 sampled sessions across the year) and
        a knowledge assessment with a pre-defined pass mark of 70/100
        administered to all participants (n~=~21 employees in the sampled
        cohort). The seven-item instrument was adapted from established
        post-training evaluation scales \citep{Kirkpatrick2010,Patel2023},
        pilot-tested with twelve technicians prior to deployment and
        achieved an internal consistency of Cronbach's $\alpha = 0.82$ in
        the present sample (above the conventional 0.70 threshold).
  \item \textbf{On-the-job observations} (Level~3): structured behaviour
        observations by trained supervisors covering the application of
        procedural knowledge during the first 90 days after each training
        intervention. Each observation used a twelve-item checklist linked
        to procedural standards. Inter-rater reliability was assessed during
        a one-day calibration exercise between two independent supervisors
        on the same set of 20 task observations and produced a Cohen's
        $\kappa = 0.78$, indicating substantial agreement.
  \item \textbf{Operational metrics} (Level~4): four metrics extracted from
        the facility's enterprise resource planning (\ERP{}) and quality
        management systems --- training failures, regulatory
        non-conformances, quality costs and turnaround time. Operational
        definitions are reported below.
  \item \textbf{Customer satisfaction surveys} (Level~4): structured surveys
        administered semi-annually to four major airline customers covering
        five categories (product delivery quality, technical documentation,
        turnaround time, customer support and engineering team assistance).
\end{enumerate}

\subsubsection*{Operational definitions of rate metrics}
\label{subsubsec:rate-definitions}

To remove ambiguity in interpretation, the rate metrics used in this study
are defined as follows:

\begin{itemize}[leftmargin=*]
  \item \textbf{Training-failure rate} = number of documented training
        failures attributable to a knowledge or skill gap (see definition
        below) divided by the total number of training-intervention exits
        in the twelve-month window, expressed per 1{,}000 maintenance tasks
        performed.
  \item \textbf{Regulatory non-conformance rate} = number of major or minor
        findings issued by the civil aviation authority (\CAA{}) during
        scheduled audit cycles in the twelve-month window, indexed to the
        number of audit cycles conducted.
  \item \textbf{First-attempt pass rate} (Level~2) = number of participants
        who reached the 70/100 pass mark on the first attempt divided by
        the total number of participants in the assessed cohort.
  \item \textbf{Work-card consultation rate} (Level~3) = proportion of
        observation events in which the technician consulted the relevant
        work card before initiating a non-routine task.
\end{itemize}

\subsubsection*{Baseline reconstruction and its limitations}
\label{subsubsec:baseline}

Baseline values for the training-failure rate and the regulatory
non-conformance rate were extracted retrospectively from the facility's
quality records using the same operational definitions adopted in the
intervention year. Baseline values for the Level~2 first-attempt pass rate
and the Level~3 work-card consultation rate were reconstructed from the
facility's \LMS{} records and from a sample of routine supervisor reports;
the latter were not collected with the structured observation protocol
used in the intervention year, and the resulting baseline figures are
therefore subject to a possible upward bias that the present study cannot
fully quantify. The reliability indices reported below (Cronbach's
$\alpha$, Cohen's $\kappa$) were calculated on the intervention-year
sample only; the baseline year did not use the same structured
instrument, so the reliability of the baseline figures is not directly
comparable.

\subsubsection*{Operational definition of \textit{training failure}}
\label{subsubsec:training-failure}

The operational definition used throughout the study is:

\begin{quote}
A \textit{training failure} is an instance where a maintenance technician,
after completing a scheduled training intervention, commits a documented
procedural error within the subsequent 90-day window that is directly
attributable to a knowledge or skill gap identified in the post-training
evaluation. Errors caused by tool malfunction, supplier non-conformance or
documented changes in the standard procedure are excluded from this
definition.
\end{quote}

Documentation of attribution is taken from the facility's root-cause analysis
records (RCA), which are independently reviewed by a quality engineer
who is not the technician's direct supervisor.

\subsection{Data analysis}
\label{subsec:analysis}

The analysis proceeded in four steps. First, a descriptive comparison of
each operational metric was performed between the baseline year and the
intervention year. Second, simple inference was applied to proportion-based
metrics: two-proportion $z$-tests (and chi-square cross-checks) with 95\%
confidence intervals were computed for the changes in training-failure
rate, non-conformance rate, first-attempt pass rate and work-card
consultation rate, to ensure that the observed differences exceed plausible
sampling variation. Third, the Kirkpatrick four-level structure was used to
organise the findings, with each level connected to specific data sources.
Fourth, an economic analysis was conducted using a standard cost-of-quality
(\CoQ{}) decomposition, presented in Section~\ref{subsec:results-economic},
to translate the operational changes into approximate annual financial and
sustainability terms.

\subsubsection*{Threats to internal validity}
\label{subsubsec:validity}

The design is a single-site pre/post comparison without a concurrent control
group; therefore, five classical threats to internal validity require
explicit consideration:
\textit{History} --- contemporaneous changes in workforce composition,
contract portfolio or supplier mix could co-vary with the intervention.
The study controls for this by holding the engine programmes and the
nominal workforce size constant within the observation window, but it
cannot fully eliminate the threat.
\textit{Maturation} --- gradual organisational learning could explain part
of the observed improvement independently of the intervention.
\textit{Instrumentation} --- changes in the way operational metrics are
recorded over time can spuriously create apparent improvement. The study
mitigates this by using metric definitions held constant from the baseline
year and audited by the facility's quality department.
\textit{Selection} --- because no random assignment is performed, the
cohort observed in the intervention year is the same that would have been
observed in the absence of the intervention; this protects against
selection bias but does not establish a counterfactual.
\textit{Hawthorne effect} --- the introduction of structured behavioural
observation may have changed technician behaviour independently of the
training content. The present study cannot disentangle the observation
effect from the training effect; future work should compare the present
observation protocol against an unobtrusive baseline (for example, an
archival document audit) to estimate the size of this contribution.
Causal attribution under these threats is necessarily provisional, and the
findings should be interpreted as \emph{associated with} the intervention
rather than \emph{produced by} it. Section~\ref{subsec:limitations} returns
to this point and indicates the quasi-experimental or
difference-in-differences extensions that would strengthen attribution in
future work.

\subsubsection*{Note on data confidentiality}

For commercial confidentiality, individual-level records of the facility
are not publicly disclosed. All values reported in the figures and tables
correspond to actual operational records of the facility, anonymised and
aggregated by the facility's quality department prior to release for
analysis. No values were synthesised or imputed.
